\section{Discussion}
\label{sec:discussion}
% Sources: false_alarm/tables/false_alarm_all.csv (mean_max_prob); detection_rules/external/*/
% misses_vs_tta8.csv (left_vs_reference, new_vs_reference = 0); compare_models/gain_difference.csv;
% pass_latency.csv; aggregate_*/uncertainty_audit_*.csv and results/paper_checks (BraTS);
% Fan et al. 2020 (random 80/10/10 image split, mDice 0.899 on CVC-ClinicDB); workshop paper
% (UNSURE 2026): larger gains claimed for the lower-performing model; UNet dDice -0.0046;
% one from-scratch model vs pretrained ones; no baseline-error vs gain analysis.

\paragraph{TTA lowered foreground probability, and the count decides how that looks.} Averaging views pulls the foreground probability down where only some views are confident. For UNet it left more images empty in nearly every checkpoint, but most of those images had not been localised before, and by the IoU criterion UNet missed fewer on ETIS-Larib and about as many elsewhere. For SwinUNETR, images with no overlap rose in all 15 checkpoints on every external set while missed images did not: that model marked a polyp on every polyp-free frame with and without TTA, so what TTA took from it was mask area and not detections. For UNet and SegFormer, whose masks could disappear, false alarms on the cecum frames fell, the same shift seen from the other side. The clearest case is within one model on one set: on ETIS-Larib the same UNet checkpoints and images give three answers, with empty predictions up from 5 to 18, no-overlap images unchanged at 33, and missed images down from 55 to 46. A study reporting only the first would conclude that TTA harms UNet. For empty outputs TTA sat close enough to the single view's detection curve that we could not separate them, so a threshold change is a plausible reading there; the lower count of missed images is not explained that way, because it depends on where the mask lands.

\paragraph{Where the effect appears.} Mostly, but not only, in models already failing. By the majority vote SegFormer changed by at most three images under each count on every set, and dropping SwinUNETR, which we broke on purpose, would leave a single two-way disagreement on ETIS-Larib. Per checkpoint the base is wider: the disagreement holds for SwinUNETR within all three collections, for UNet across two, and SegFormer's empty predictions rise on two before correction. The effect is still largest where the model is weakest, which narrows the claim without removing it, because a pre-deployment check is run precisely where a model may be weak or shifted.

\paragraph{Clinical reading.} Both errors have a clinical cost. A higher adenoma detection rate is linked to less interval cancer \cite{corley2014adr}, but that evidence is per endoscopist and not per frame, so our data cannot weigh one error against the other. For UNet on ETIS-Larib, TTA left about 66 more images per 1000 empty while missing about 46 fewer; the false alarms it removed were counted on other frames in another hospital, so those are separate observations and not one trade. A frame on which the system goes quiet still gives no prompt, so empty predictions deserve their own count even when the removed prediction was wrong. Eight views took about 298\,ms per frame at batch 1, so in this form TTA suits offline use rather than live video, though batching could change that and we did not measure it. Two points follow. Report misses with an explicit localisation criterion, next to empty predictions and false alarms, whenever the inference procedure changes, since any single rule would have settled the verdict here by itself. And check empty predictions separately, because no ROI signal can score them.

\paragraph{Limitations.} This is a retrospective study on public data, with no reader study and no patient-level outcome. The central negative result is a failure to detect harm, not a demonstration of safety: no non-inferiority margin, no power for one, and intervals still admitting about 20 more missed images per 1000. The 15 checkpoints are 5 folds $\times$ 3 seeds from one recipe on one collection and not independent replicates, so every test here describes this pool. We count misses per image, so a missed frame is not a missed polyp, and image-level IoU is not a clinical standard, is harder to reach for small polyps and ignores images with several polyps; we added that criterion after seeing the any-overlap counts and swept it on the data that motivated it. A false alarm is any non-empty mask, the cecum frames carry an anatomical-landmark label from one hospital, and the ET-free population is 33 patients. All 2D models come from one small single-centre collection, no external collection carries centre or sequence labels, and the external evidence rests on well under 100 patients; PolypGen \cite{ali2023polypgen} would permit a cluster-correct, multi-centre analysis. We did not analyse calibration or re-threshold non-empty predictions, so missed images are compared at matched thresholds and not at matched false-alarm rates, and part of the external distance gain may be survivorship. The pattern rests on three models, one failed by design; the uncertainty audit is internal only; and we studied geometric TTA only.
